% ------------------------------ 76 characters -----------------------------
\documentclass{article}

% ------------------------------ 76 characters -----------------------------
\usepackage[l2tabu, orthodox]{nag}

% -- Texto e codificação
\usepackage[utf8]{inputenc}
\usepackage[portuguese]{babel}
\usepackage{microtype}
\usepackage{xcolor}
\usepackage{anyfontsize}

% -- Tipo de letra
\usepackage{heuristica}
\usepackage[heuristica,vvarbb,bigdelims]{newtxmath}
\usepackage[T1]{fontenc}
\renewcommand*\oldstylenums[1]{\textosf{#1}}

%\usepackage{lmodern}
%\usepackage[T1]{fontenc}

% -- Definir margens do documento
\usepackage{geometry}
\geometry{verbose, nomarginpar,
    tmargin = 2.5cm,
    bmargin = 2.5cm,
    lmargin = 2.5cm,
    rmargin = 2.5cm}
%\usepackage{showframe}

% -- Cabeçalho e rodapé
\usepackage{fancyhdr}
\fancyhf{}
\fancyhead[L]{\textsc{Instituto Superior Técnico}}
\fancyhead[R]{\textsc{Programação de Sistemas}}
\fancyfoot[C]{\thepage}
\pagestyle{fancy}

% -- Funções matemáticas extra
\usepackage{mathtools}
\usepackage{siunitx}

% -- Símbolos extra
\usepackage{amssymb}
\usepackage{textcomp}
\usepackage{gensymb}
\usepackage{cancel}

% -- Referências
\usepackage{hyperref}

% --  Definições de imagens
\usepackage{graphicx}
\graphicspath{{graphics/}}
\usepackage{caption}
\usepackage{subcaption}

% -- Desenhar circuitos elétricos e lógicos
\usepackage{tikz}
\usepackage{pgfplots}
\usetikzlibrary{arrows.meta,positioning}
\pgfplotsset{compat=1.5}
\pgfplotsset{table/search path = {data}}
\pgfplotsset{	/pgf/number format/use comma,}

% -- Integrar código fonte
\usepackage{minted}


\begin{document}

\thispagestyle{empty}
	\includegraphics[viewport=9.5cm 11cm 0cm 0cm,scale=0.29]{IST_A_CMYK_POS}
	
	\begin{center}
		\vspace{70mm} % --  Espaço em branco
		\rule{\linewidth}{0.5pt} \\
		\Huge \textsc{Clipboard Distribuído} \\
		\rule{\linewidth}{2pt} \\
		\vspace{10mm} % -- Espaço em branco
		\LARGE Projecto Final
		
		\vspace{\fill} % --  Espaço em branco variável
		\large
		Daniel de Schiffart \texttt{81479} \\
		João Gonçalves \texttt{81040}
		
		\vspace{10mm} % --  Espaço em branco
		\Large Instituto Superior Técnico \\
		\textit{Mestrado em Engenharia Aeroespacial} \\
		\vspace{1mm}
		\large Programação de Sistemas
		
		\vspace{10mm} % --  Espaço em branco
		\Large Maio de 2018
	\end{center}
\newpage

\section{Objetivo}

Para o projeto final da unidade curricular de Programação de Sistemas do
segundo semestre do ano curricular de 2017/2018 foi desenvolvido e
implementado um \textit{clipboard} distribuído em C. O objetivo do
projeto proposto consistia na implementação do ante-referido 
\textit{clipboard} através de uma biblioteca, uma API e um processo
local que controlassem o funcionamento do mesmo.

O funcionamento do \textit{clipboard} distribuído é análogo ao
funcionamento da função de copiar-colar implementada em muitos diferentes
programas e sistemas. Uma API permite a outras aplicações comunicar com
um processo local (o \textit{clipboard} local) e guardar e retirar dados
temporariamente, que serão sincronizados com outras instâncias do mesmo
\textit{clipboard} local que estão em funcionamento em outros sistemas,
desde que estejam conectados entre si. O conjunto das instâncias deste
\textit{clipboard} terá então os mesmos dados guardados em todos os
sistemas, facilitando assim a comunicação de dados entre estes.
Cada conjunto de \textit{clipboards} interligados tem espaço para
armazenamento de dez variáveis diferentes, identificados de 1 a 10.

\end{document}